% =====================================================================
%  Desenho da base de dados - COP30 e emprego em Alojamento e Alimentação
%  Compilar no Overleaf com pdfLaTeX (padrão)
% =====================================================================
\documentclass[11pt,a4paper]{article}

\usepackage[utf8]{inputenc}
\usepackage[T1]{fontenc}
\usepackage[brazil]{babel}
\usepackage{lmodern}
\usepackage{microtype}
\usepackage[top=2.5cm, bottom=2.5cm, left=2.5cm, right=2.5cm]{geometry}
\usepackage{setspace}
\setstretch{1.15}
\usepackage{amsmath, amssymb}
\usepackage{booktabs, tabularx, array}
\usepackage{float}
\usepackage{enumitem}
\usepackage{xcolor}
\usepackage{tikz}
\usetikzlibrary{positioning, arrows.meta}
\usepackage{titlesec}
\usepackage[hidelinks]{hyperref}

% ---- Cores e estilos ---------------------------------------------------
\definecolor{azul}{HTML}{1F5FAD}
\definecolor{azulclaro}{HTML}{E6EEF8}
\definecolor{cinza}{HTML}{6B6B6B}
\definecolor{borda}{HTML}{B9B6AC}
\definecolor{verde}{HTML}{2E8B3D}
\definecolor{amarelo}{HTML}{E0A000}

\titleformat{\section}{\large\bfseries\color{azul}}{\thesection}{0.6em}{}
\titleformat{\subsection}{\normalsize\bfseries}{\thesubsection}{0.6em}{}
\setlist{itemsep=2pt, topsep=4pt}
\newcolumntype{L}{>{\raggedright\arraybackslash}X}
\newcommand{\var}[1]{\texttt{#1}}
\newcommand{\pronto}[1]{\par\smallskip\noindent\textcolor{verde}{\textbf{Pronto quando:}} #1}
\renewcommand{\arraystretch}{1.25}

\tikzset{
  caixa/.style={draw=borda, rounded corners=4pt, line width=0.8pt, align=left,
                text width=3.2cm, minimum height=2.3cm, inner sep=5pt, font=\footnotesize,
                execute at begin node={\hyphenpenalty=10000\exhyphenpenalty=10000}},
  destaque/.style={caixa, draw=azul, fill=azulclaro, line width=1.4pt},
  seta/.style={-{Latex[length=2.2mm]}, draw=cinza, line width=0.8pt}
}
\newcommand{\nabase}{\textcolor{verde}{\large$\bullet$}\,na base}
\newcommand{\acoletar}{\textcolor{amarelo}{\large$\bullet$}\,a coletar}

% =====================================================================
\begin{document}

\begin{center}
  {\LARGE\bfseries\color{azul} Desenho da base de dados}\\[4pt]
  {\large COP30 e emprego formal em Alojamento e Alimentação}\\[8pt]
  Marcela Paz da Silva Enriquez \quad·\quad Outubro de 2026
\end{center}

\vspace{6pt}
\noindent\fcolorbox{azul}{azulclaro}{\parbox{\dimexpr\linewidth-2\fboxsep-2\fboxrule}{
\textbf{Em uma frase:} a base é um painel mensal das 27 capitais brasileiras (26 estaduais e Brasília),
de janeiro de 2021 a julho de 2026, que compara o emprego formal em Alojamento e Alimentação de Belém
com o de cidades que não sediaram a COP30.}}

% ---------------------------------------------------------------------
\section{Objetivo e pergunta}

\textbf{Pergunta de pesquisa:} o fim da COP30 esteve associado a uma redução do saldo mensal de
empregos formais em Alojamento e Alimentação em Belém e no Pará? Em que magnitude?

Para responder, a base precisa permitir três coisas:
\begin{enumerate}
  \item \textbf{Medir o resultado} com a frequência certa: emprego formal mensal, por setor e por
        município, porque efeitos de megaeventos são curtos e setoriais (Feddersen e Maennig, 2013).
  \item \textbf{Construir um contrafactual:} o que teria acontecido com Belém sem a COP, a partir de
        outras capitais (Hagn e Maennig, 2009; Bronzini, Mocetti e Mongardini, 2019).
  \item \textbf{Mostrar o mecanismo:} os elos entre a chegada dos visitantes e as contratações.
\end{enumerate}

% ---------------------------------------------------------------------
\section{Estrutura da base}

A unidade de observação é capital $\times$ mês: 27 capitais, de jan/2021 a jul/2026, com a análise
a partir de jan/2022.

\begin{figure}[H]
\centering
\begin{tikzpicture}[node distance=0.4cm]
  \node[caixa] (f) {\textbf{1. Fontes}\\[2pt] CAGED: MOV, FOR, EXC\\ PMS/IBGE, tab.\ 5906\\[2pt]
        \textcolor{cinza}{A coletar: RAIS, ANAC, IPCA, SICONFI}};
  \node[caixa, right=of f] (l) {\textbf{2. Limpeza}\\[2pt] Capitais, CNAE 55+56\\ FOR soma, EXC subtrai\\
        Pelo mês do movimento\\[2pt] \textcolor{cinza}{5 checagens}};
  \node[destaque, right=of l] (p) {\textbf{3. Painel}\\[2pt] 27 capitais\\ $\times$ 67 meses\\
        Saldo normalizado\\ Dummies de fase};
  \node[caixa, right=of p] (a) {\textbf{4. Análises}\\[2pt] DiD por fase\\ Estudo de evento\\
        Controle sintético\\ Placebos (p-valor)};
  \draw[seta] (f) -- (l); \draw[seta] (l) -- (p); \draw[seta] (p) -- (a);
\end{tikzpicture}
\caption{Construção da base em quatro etapas.}
\end{figure}

Os microdados chegam como um registro por contratação ou desligamento no país inteiro; a limpeza
filtra capitais e setores, corrige as declarações atrasadas e entrega um painel balanceado.

% ---------------------------------------------------------------------
\section{Mecanismo e blocos de variáveis}

Cada elo do caminho entre a COP30 e o emprego tem um indicador. Hoje a base mede o último elo e o
controle de demanda; os elos intermediários são a próxima coleta.

\begin{figure}[H]
\centering
\begin{tikzpicture}[node distance=0.4cm]
  \node[font=\footnotesize\bfseries, text=cinza, anchor=west] at (-1.75,1.55) {Canal da demanda turística};
  \node[caixa] (v) {\textbf{Visitantes chegam}\\[2pt] Passageiros no aeroporto (ANAC)\\[3pt] \acoletar};
  \node[caixa, right=of v] (d) {\textbf{Demanda sobe}\\[2pt] Ocupação hoteleira; IPCA hospedagem\\[3pt] \acoletar};
  \node[caixa, right=of d] (e) {\textbf{Empresas vendem}\\[2pt] Volume de serviços (PMS/IBGE)\\[3pt] \nabase};
  \node[destaque, right=of e] (c) {\textbf{Empresas contratam}\\[2pt] Admissões, desligamentos e saldo (CAGED)\\[3pt] \nabase};
  \draw[seta] (v) -- (d); \draw[seta] (d) -- (e); \draw[seta] (e) -- (c);

  \node[font=\footnotesize\bfseries, text=cinza, anchor=west] at (-1.75,-1.75) {Canal do setor público};
  \node[caixa, below=1.1cm of v] (g) {\textbf{Governo investe}\\[2pt] Estado e prefeitura (SICONFI)\\[3pt] \acoletar};
  \node[caixa, right=of g] (o) {\textbf{Obras}\\[2pt] Infraestrutura para o evento};
  \node[caixa, right=of o] (k) {\textbf{Emprego em obras}\\[2pt] Construção civil, CNAE 41 a 43\\[3pt] \acoletar};
  \node[caixa, right=of k] (lg) {\textbf{Legado?}\\[2pt] Uso das obras, salários e setores};
  \draw[seta] (g) -- (o); \draw[seta] (o) -- (k); \draw[seta] (k) -- (lg);
\end{tikzpicture}
\caption{A COP30 chega ao emprego por dois canais. Depois do evento, o canal da demanda corre ao
contrário: visitantes saem, a demanda cai e as contratações desaceleram.}
\end{figure}

% ---------------------------------------------------------------------
\section{Fontes de dados}

\begin{table}[H]
\centering\small
\begin{tabularx}{\linewidth}{>{\raggedright}p{3.1cm}LLp{2.3cm}l}
\toprule
\textbf{Fonte} & \textbf{O que traz} & \textbf{Papel} & \textbf{Recorte} & \textbf{Situação}\\
\midrule
Novo CAGED (microdados) & Cada admissão e desligamento formal, com CNAE, município e salário & Resultado; emprego em obras & Vínculo, mensal & Na base\\
PMS/IBGE, tab.\ 5906 & Índice de volume de serviços (2022 = 100) & Controle de demanda & UF, mensal & Na base\\
RAIS (PDET/MTE) & Estoque de vínculos, salários e ocupações (CBO) & Normalização; dimensões do trabalho & Vínculo, anual & A coletar\\
ANAC & Passageiros por aeroporto & Elo: chegada de visitantes & Aeroporto, mensal & A coletar\\
IPCA/IBGE & Preços de hospedagem e alimentação fora do domicílio & Elo: pressão sobre preços & RM, mensal & A coletar\\
SETUR-PA e redes hoteleiras & Ocupação hoteleira & Elo: demanda por hospedagem & Belém, mensal & A pedir\\
SICONFI/Tesouro & Investimento do estado e do município & Canal público & Ente, bimestral & A coletar\\
PNAD Contínua; DIEESE-PA & Informalidade no setor & Descrição do que o CAGED não vê & UF, trimestral & A coletar\\
\bottomrule
\end{tabularx}
\end{table}

O CAGED é usado pelos microdados porque permite recortar por subclasse e município e porque soma as
declarações fora do prazo e desconta as exclusões pelo mês do movimento (Portaria MTE nº 2.420/2023).
A RAIS não substitui o CAGED: é anual e completa, enquanto o CAGED é mensal; a RAIS entra como
estoque inicial e para salários e ocupações.

% ---------------------------------------------------------------------
\section{Dicionário de variáveis}

\begin{table}[H]
\centering\small
\begin{tabularx}{\linewidth}{p{3.3cm}p{2.1cm}Lp{2cm}}
\toprule
\textbf{Variável} & \textbf{Bloco} & \textbf{Definição} & \textbf{Fonte}\\
\midrule
\var{cap}, \var{mun} & Identificação & Nome e código IBGE (6 dígitos) da capital & IBGE\\
\var{competencia}, \var{mes\_ano} & Identificação & Mês de competência e mês do ano (1 a 12) & ---\\
\var{adm}, \var{desl} & Resultado & Admissões e desligamentos em CNAE 55 + 56 & CAGED\\
\var{saldo} & Resultado & $adm - desl$ & CAGED\\
\var{adm\_ref} & Resultado & Média mensal de admissões, jan/2022 a jun/2024 & CAGED\\
\var{saldo\_norm} & Resultado & $100 \times saldo / adm\_ref$: \% de um mês normal de contratações & CAGED\\
\var{emp\_idx} & Resultado & Saldo acumulado desde jan/2022, normalizado & CAGED\\
\var{saldo\_55}, \var{saldo\_56} & Robustez & Saldo de Alojamento e de Alimentação & CAGED\\
\var{estoque} & Normalização & Vínculos ativos em 31/12 do ano anterior & RAIS\\
\var{sal\_adm} & Mecanismo & Salário médio de admissão & CAGED\\
\var{saldo\_constr} & Mecanismo & Saldo na construção civil (CNAE 41 a 43) & CAGED\\
\var{pax} & Mecanismo & Passageiros no aeroporto da capital & ANAC\\
\var{ocup\_hotel} & Mecanismo & Taxa de ocupação hoteleira & SETUR-PA\\
\var{ipca\_hosp} & Mecanismo & Variação mensal dos preços de hospedagem & IPCA\\
\var{share55} & Controle & Participação de Alojamento nas admissões & CAGED\\
\var{ln\_ivs} & Controle & Log do índice de volume de serviços da UF & PMS/IBGE\\
\var{tratado} & Tratamento & 1 para Belém, 0 para as demais & ---\\
\var{D\_prep}, \var{D\_evento}, \var{D\_desmob} & Tratamento & Fases da COP30 (seção 6) & ---\\
\bottomrule
\end{tabularx}
\end{table}

A normalização de \var{saldo\_norm} torna as capitais comparáveis: $+300$ vagas é muito em Palmas e
pouco em São Paulo, mas ``$+20\%$ de um mês normal'' significa o mesmo nas duas.

% ---------------------------------------------------------------------
\section{Tratamento, fases e grupo de controle}

Belém é a única unidade tratada; as outras 25 capitais (todas menos o Rio de Janeiro) formam o grupo
de controle.

\begin{table}[H]
\centering\small
\begin{tabularx}{\linewidth}{p{3cm}p{3.6cm}L}
\toprule
\textbf{Fase} & \textbf{Período} & \textbf{Uso}\\
\midrule
Referência & jan/2022 a jun/2024 & Linha de base; normaliza o saldo e ajusta o controle sintético\\
Preparação & jul/2024 a out/2025 & Antecipação: contratações pela demanda esperada\\
Evento & nov/2025 & Mês da conferência\\
Desmobilização & dez/2025 a jul/2026 & Fase de interesse da pergunta\\
\bottomrule
\end{tabularx}
\end{table}

\begin{itemize}
  \item \textbf{2021 fica fora da análise:} a recuperação da pandemia foi desigual entre capitais.
  \item \textbf{Rio de Janeiro sai do controle:} sediou a cúpula do G20 em novembro de 2024.
  \item \textbf{A janela de preparação será testada:} o aquecimento parece concentrado em set--nov/2025.
\end{itemize}

% ---------------------------------------------------------------------
\section{Estratégia de identificação}

\textbf{Diferenças em diferenças por fase} (capital $c$, mês $t$):
\begin{equation}
y_{ct} = \alpha_{c,\,m(t)} + \lambda_t + \beta_P\,(B_c \times Prep_t) + \beta_E\,(B_c \times Evento_t)
       + \beta_D\,(B_c \times Desmob_t) + \varepsilon_{ct}
\end{equation}
$\alpha_{c,m(t)}$ é um efeito fixo de capital $\times$ mês do ano (sazonalidade de cada cidade: Círio em
Belém, Carnaval em Salvador); $\lambda_t$ é um efeito fixo de mês calendário (choques comuns); $B_c = 1$
para Belém. O coeficiente de interesse é $\beta_D$.

\textbf{Estudo de evento:} uma dummy para cada mês de Belém a partir de jul/2024.
\textbf{Controle sintético:} uma Belém artificial, média ponderada das capitais que melhor reproduzem
Belém de jan/2022 a jun/2024 (variável \var{emp\_idx}).

\begin{table}[H]
\centering\small
\begin{tabularx}{\linewidth}{p{3.2cm}LL}
\toprule
\textbf{Hipótese} & \textbf{O que exige} & \textbf{Como é verificada}\\
\midrule
Tendências paralelas & Sem a COP, Belém seguiria as outras capitais & Coeficientes do estudo de evento perto de zero antes de jul/2024\\
Sem antecipação fora da janela & Nenhum efeito antes da preparação & Mesmo teste, com a preparação antecipada\\
Sem contaminação & Controles não afetados pela COP & Rio excluído; reestimar sem cada doadora relevante\\
\bottomrule
\end{tabularx}
\end{table}

\textbf{Inferência:} com uma única unidade tratada, o erro-padrão agrupado não é confiável (Conley e
Taber, 2011). O p-valor vem de placebos: o modelo é reestimado fingindo que cada capital sediou a COP.
Com 26 unidades, ficar em 1º lugar corresponde a $p \approx 0{,}04$.

% ---------------------------------------------------------------------
\section{Leitura macroeconômica}

\begin{itemize}
  \item \textbf{Demanda derivada:} o emprego segue a demanda por diárias e refeições (Hicks, 1932); o
        índice de serviços capta o ciclo regional e os efeitos fixos de mês, os choques nacionais.
  \item \textbf{Expectativas:} contratações pela demanda esperada explicam a antecipação; o emprego se
        ajusta com defasagem (Alemayehu e Tveteraas, 2020).
  \item \textbf{Fluxos e ciclo:} a rotatividade na hospitalidade é a mais sensível ao ciclo (Dogru et
        al., 2023).
  \item \textbf{Agregação:} efeitos locais se diluem em unidades grandes (Hagn e Maennig, 2009).
  \item \textbf{Composição setorial:} no Jubileu de 2000, o emprego de Roma cresceu em setores de baixa
        produtividade (Bronzini et al., 2019); salário de admissão e construção testam isso em Belém.
\end{itemize}

% ---------------------------------------------------------------------
\section{Checagens de qualidade e limitações}

\begin{itemize}[label=$\square$]
  \item Nenhuma capital com admissões zeradas (erro de código IBGE).
  \item Painel balanceado: 27 capitais $\times$ todos os meses com arquivo MOV.
  \item Meses sem arquivo MOV ficam ausentes, nunca zero.
  \item Alojamento + Alimentação = total em todas as linhas.
  \item Saldo de Belém igual ao da base anterior nos mesmos meses.
\end{itemize}

\begin{table}[H]
\centering\small
\begin{tabularx}{\linewidth}{p{3.2cm}LL}
\toprule
\textbf{Limitação} & \textbf{Consequência} & \textbf{Tratamento}\\
\midrule
Só emprego formal & Informalidade fica fora & Descrita com PNAD Contínua e DIEESE-PA\\
8 meses pós-evento & Baixo poder estatístico & Não significância $\neq$ ausência de efeito\\
Uma unidade tratada & Erro-padrão convencional inválido & Inferência por placebos\\
Revisões do CAGED & Meses recentes ainda mudam & Reprocessar a cada divulgação\\
\bottomrule
\end{tabularx}
\end{table}

% ---------------------------------------------------------------------
\section{Contexto e evidências complementares}

\textbf{Métodos ex ante do turismo.} Insumo-produto (IO), equilíbrio geral computável (CGE), matriz de
contabilidade social (SAM), conta satélite do turismo (TSA), pesquisas de gasto e multiplicadores de
valor adicionado (GVA). Partem do gasto turístico e projetam efeitos; esta pesquisa é ex post e os usa
como contraponto.

\textbf{Dimensões do trabalho.} Emprego, salário e renda, produtividade, estrutura ocupacional,
desigualdade de ganhos, características da força de trabalho e condições de trabalho.

\textbf{Eventos de comparação.}
\begin{table}[H]
\centering\small
\begin{tabularx}{\linewidth}{p{3.2cm}LL}
\toprule
\textbf{Evento} & \textbf{O que ensina} & \textbf{Uso na pesquisa}\\
\midrule
Círio de Nazaré & Grande evento recorrente; a COP não se repete & Sazonalidade de outubro, absorvida pelos efeitos fixos\\
Carnaval e Parintins & Aumento de renda concentrado em poucos dias & Comparação: evento recorrente $\times$ evento único\\
Rio 2016 e reformas & Investimento concentrado e deslocamento de moradores & Literatura de legado; Rio fora do controle\\
\bottomrule
\end{tabularx}
\end{table}

\textbf{Infraestrutura e legado.} Além das receitas, há obras, construtoras e novos negócios. Ficam em
aberto o uso das construções, os impactos ambientais e culturais, o efeito simbólico sobre o imaginário
da cidade e a justificativa do investimento público.

\textbf{Contatos:} DIEESE-PA, SETUR-PA, redes hoteleiras e Paula (IPEA), da palestra do PET sobre
megaeventos.

% ---------------------------------------------------------------------
\section{Passo a passo para os próximos passos}

\begin{enumerate}[label=\textbf{\arabic*.}, leftmargin=*, itemsep=6pt]
\item \textbf{Organizar a pasta.} Tudo em \texttt{Documentos/cop30} (sem espaço), que guarda os dados do CAGED.
  \begin{itemize}
    \item Mova os PDFs e textos da pasta \texttt{cop 30} (com espaço) para \texttt{cop30}; nunca o contrário.
    \item Crie as subpastas \texttt{scripts}, \texttt{textos} e \texttt{saidas}.
  \end{itemize}
  \pronto{só existe a pasta \texttt{cop30}, com \texttt{dados\_caged} dentro.}

\item \textbf{Rodar o painel das capitais.} Abra \texttt{03\_belem\_vs\_capitais.R} e clique em
  \textbf{Source} (30 a 60 minutos). Na tabela inicial, nenhuma capital pode ter zero admissões.
  \pronto{\texttt{list.files(pattern = "\^{}did\_|\^{}sc\_")} lista os arquivos de resultado.}

\item \textbf{Conferir a base.} Rode as cinco checagens da seção 9.
  \pronto{as cinco caixas estão marcadas.}

\item \textbf{Buscar a RAIS.} No PDET/MTE, vínculos ativos em 31/12/2020 e 31/12/2021 em CNAE 55 + 56
  por capital, com salário médio e ocupação (CBO). O estoque de Belém vai para \texttt{ESTOQUE\_INICIAL}.
  \pronto{existe \texttt{rais\_capitais.csv} com capital, ano, vínculos e salário.}

\item \textbf{Coletar os elos do mecanismo.}
  \begin{itemize}
    \item Passageiros: ANAC, dados estatísticos do transporte aéreo, por aeroporto e mês.
    \item Preços: IPCA de Belém no SIDRA, subitens de hospedagem e alimentação fora do domicílio
          (confirmar o número da tabela).
    \item Ocupação hoteleira: séries mensais da SETUR-PA e das redes hoteleiras.
    \item Investimento público: SICONFI, estado do Pará e prefeitura de Belém.
  \end{itemize}
  \pronto{cada fonte vira um CSV com \texttt{competencia} e a variável, de jan/2022 a jul/2026.}

\item \textbf{Descrever a informalidade.} PNAD Contínua (trimestral) e DIEESE-PA: informalidade em
  alojamento e alimentação no Pará.
  \pronto{um parágrafo e um gráfico mostram o tamanho do setor que o CAGED não vê.}

\item \textbf{Estimar, testar e escrever.}
  \begin{itemize}
    \item DiD por fase, estudo de evento e controle sintético (script 03).
    \item Robustez: preparação curta (set--out/2025), sem cada doadora com peso alto, só alojamento.
    \item Levar aos professores: janela de preparação, regiões metropolitanas e capitais excluídas.
  \end{itemize}
  \pronto{os resultados entram no capítulo de resultados e nos slides.}
\end{enumerate}

% ---------------------------------------------------------------------
\section*{Referências}
\begin{small}
\setlength{\parindent}{0pt}\setlength{\parskip}{4pt}\raggedright

ALEMAYEHU, F. K.; TVETERAAS, S. L. Long-run labour flexibility in hospitality: a dynamic common correlated effects approach. \textbf{Tourism Economics}, v. 26, n. 4, p. 704-718, 2020.

BRASIL. Ministério do Trabalho e Emprego. \textbf{Portaria MTE nº 2.420, de 10 de julho de 2023}. Brasília, 2023.

BRONZINI, R.; MOCETTI, S.; MONGARDINI, M. \textbf{The economic effects of big events}: evidence from the Great Jubilee 2000 in Rome. Roma: Banca d'Italia, 2019. (Temi di discussione, n. 1208).

CONLEY, T. G.; TABER, C. R. Inference with ``difference in differences'' with a small number of policy changes. \textbf{Review of Economics and Statistics}, v. 93, n. 1, p. 113-125, 2011.

DOGRU, T.; McGINLEY, S.; SHARMA, A.; IŞIK, C.; HANKS, L. Employee turnover dynamics in the hospitality industry vs. the overall economy. \textbf{Tourism Management}, v. 99, 104783, 2023.

FEDDERSEN, A.; MAENNIG, W. Mega-events and sectoral employment: the case of the 1996 Olympic Games. \textbf{Contemporary Economic Policy}, v. 31, n. 3, p. 580-603, 2013.

HAGN, F.; MAENNIG, W. Large sport events and unemployment: the case of the 2006 soccer World Cup in Germany. \textbf{Applied Economics}, v. 41, n. 25, p. 3295-3302, 2009.

HICKS, J. R. \textbf{The theory of wages}. London: Macmillan, 1932.
\end{small}

\end{document}
